\section{Выбор средств реализации}

На основе анализа требований и существующих технологий были выбраны следующие инструменты для разработки библиотеки \texttt{cast-castle}. Каждый инструмент был тщательно оценён с точки зрения производительности, совместимости с Kotlin и простоты использования в задачах метапрограммирования.

\subsection{Kotlin Symbol Processing (KSP)}

KSP выбран в качестве основного инструмента для анализа кода и генерации файлов \cite{refg, refj}. В отличие от предыдущего поколения инструментов (KAPT), KSP работает напрямую с Kotlin-символами, такими как \texttt{KSClassDeclaration}, \texttt{KSFunctionDeclaration}, \texttt{KSPropertyDeclaration}, не выполняя преобразования в Java AST. Это даёт несколько ключевых преимуществ для метапрограммирования:

\begin{itemize}
    \item \textbf{Скорость}. Поскольку нет этапа конвертации Kotlin-кода в Java-модели, обработка символов происходит значительно быстрее. В крупных проектах разница может составлять десятки процентов времени компиляции.
    \item \textbf{Сохранение Kotlin-семантики}. KSP «видит» все особенности Kotlin: nullable-типы (\texttt{String?}), не-шаблонные классы, inline-классы, sealed-классы, функции расширения, параметры по умолчанию и т.д. Это критически важно для правильного анализа полей классов и генерации корректного кода мапперов \cite{refa}.
    \item \textbf{Incremental processing}. KSP поддерживает инкрементальную компиляцию: при изменении одного исходного файла переобрабатываются только те символы, которые реально зависят от изменений. Это ускоряет итерации разработки.
    \item \textbf{Простота API}. Интерфейсы KSP интуитивно понятны для тех, кто знаком с Kotlin-рефлексией. Например, чтобы получить все функции класса, достаточно вызвать \texttt{getDeclaredFunctions()}.
\end{itemize}

Библиотека \texttt{cast-castle} использует KSP версии \texttt{2.0.0-1.0.21}, совместимой с Kotlin 2.0.0. KSP подключается через Gradle-плагин \texttt{com.google.devtools.ksp}.

\subsection{KotlinPoet}

KotlinPoet — это библиотека с открытым исходным кодом, разработанная компанией Square для генерации исходных файлов Kotlin \cite{refg, refk}. Она задумывалась как Kotlin-аналог популярной библиотеки JavaPoet и предоставляет типобезопасный API для программного создания корректного и читаемого кода. В контексте \texttt{cast-castle} KotlinPoet используется для генерации вспомогательных функций-мапперов и классов-реализаций.

\subsubsection{Основные компоненты KotlinPoet}

В основе библиотеки лежат абстракции, называемые «Spec»-классами, каждая из которых отвечает за определённую конструкцию языка Kotlin:

\begin{itemize}
    \item \textbf{FileSpec} — представляет целый Kotlin-файл (`.kt`). Он служит корневым контейнером для пакета, импортов и всех объявлений типов, функций и свойств верхнего уровня.
    \item \textbf{TypeSpec} — описывает тип: класс, интерфейс, объект, enum или аннотацию. Позволяет задавать модификаторы (например, \texttt{abstract}, \texttt{sealed}, \texttt{data}), первичный конструктор, наследование и внутренних членов (функций и свойств).
    \item \textbf{FunSpec} — генерирует функции, включая конструкторы, геттеры и сеттеры. Позволяет задать имя, модификаторы, параметры (через \texttt{ParameterSpec}), возвращаемый тип, тело функции (в виде \texttt{CodeBlock}) и аннотации.
    \item \textbf{PropertySpec} — генерирует свойства (\texttt{val} или \texttt{var}) с инициализатором, геттером и сеттером, код которых также задаётся через \texttt{CodeBlock}.
    \item \textbf{ParameterSpec} — описывает параметры функций и конструкторов: имя, тип, значение по умолчанию, модификатор \texttt{vararg}.
    \item \textbf{CodeBlock} — ключевой элемент для построения тела функций. Использует специальную систему форматирования с заполнителями: \texttt{\%L} для литералов, \texttt{\%N} для имён, \texttt{\%S} для строк, \texttt{\%T} для типов. Для создания конструкций с отступами и скобками предоставляются методы \texttt{beginControlFlow()} и \texttt{endControlFlow()}.
\end{itemize}

\subsubsection{Процесс генерации кода}

Работа с KotlinPoet напоминает строительство иерархической структуры: от конкретных частей к файлу целиком. Сначала создаются \texttt{FunSpec} и \texttt{PropertySpec}, затем из них собирается \texttt{TypeSpec}, который помещается в \texttt{FileSpec}. Финальный код генерируется вызовом \texttt{writeTo()} для записи в файл или \texttt{toString()} для получения строки.

\subsubsection{Интеграция с KSP}

KotlinPoet идеально сочетается с KSP: модуль \texttt{kotlinpoet-ksp} предоставляет утилиты для преобразования KSP-типов (\texttt{KSType}) в типы KotlinPoet (\texttt{ClassName}, \texttt{ParameterizedTypeName}) и для записи сгенерированных файлов напрямую в \texttt{CodeGenerator} KSP. Это позволяет избежать ручного управления файлами и автоматически учитывать зависимости между исходными и сгенерированными файлами.

\subsection{Gradle и многомодульная структура}

Для сборки и распространения библиотеки используется система сборки Gradle \cite{refm}. Проект разделён на два независимых модуля:

\begin{itemize}
    \item \texttt{:cast-castle-annotations} — содержит только аннотацию \texttt{@CastCastleMapper} и не имеет зависимостей. Этот модуль подключается к проекту пользователя как обычная \texttt{implementation} зависимость. Он должен присутствовать как на этапе компиляции, так и в рантайме.
    \item \texttt{:cast-castle-processor} — содержит KSP-процессор, все вспомогательные классы (модели, резолверы, генератор, утилиты) и зависимость от модуля \texttt{:annotations}. Подключается к проекту пользователя через специальную конфигурацию \texttt{ksp}, предоставляемую плагином \texttt{com.google.devtools.ksp}. Такое разделение гарантирует, что код процессора не попадёт в итоговый APK или JAR рантайма, поскольку \texttt{ksp} зависимость обрабатывается только на этапе компиляции.
\end{itemize}

Такая архитектура является стандартной для KSP-библиотек и рекомендуется Google \cite{refc}. Она обеспечивает чистоту зависимостей и минимизирует размер итогового приложения. Применение шаблонов проектирования, в частности фабричного метода для создания процессора, обусловлено требованиями гибкости и расширяемости \cite{refe}.

\subsection{Стандартная библиотека Kotlin}

Для выполнения базовых операций (работа с коллекциями, строками, null-безопасность) используется стандартная библиотека Kotlin (stdlib) \cite{refa}. Она подключается автоматически через Gradle и не требует дополнительных действий. В процессоре активно используются функции работы с коллекциями (\texttt{filter}, \texttt{map}, \texttt{groupBy}) и строковые операции (построение имён, форматирование). Стандартная библиотека Kotlin также предоставляет корутины для асинхронной обработки.

\subsection{Итоговый стек технологий}

Суммируя, стек технологий библиотеки \texttt{cast-castle} выглядит следующим образом:

\begin{itemize}
    \item Язык реализации: Kotlin \cite{refa}
    \item Метапрограммирование: KSP \cite{refg, refj}
    \item Генерация кода: KotlinPoet \cite{refk}
    \item Сборка: Gradle \cite{refm} с плагином \texttt{com.google.devtools.ksp}
    \item Тестирование: \texttt{ksp-compiler-testing} и примеры в папке \texttt{samples}
\end{itemize}

Все выбранные инструменты являются актуальными, активно поддерживаются сообществом и широко используются в современных Kotlin-проектах. Это обеспечивает долгосрочную поддержку библиотеки и её совместимость с будущими версиями экосистемы.
